% Theorem-like environments, the "problem" box and numbering.
% Input after cleveref: theorem definitions must follow cleveref, otherwise every environment that
% shares a counter is referred to by the same name.
% \declaretheorem (thmtools), not \newtheorem: with the LaTeX 2025-11-01 kernel and cleveref 0.21.4,
% \cref names a lemma that shares the definition counter "Definition" (tested 2026-10-07);
% sibling= gives each environment its own name and plural ("Lemmas 2.3 and 2.4").
% Styles are the theorem styles of the class sn-jnl, assigned as in its sample sn-article.tex:
% thmstyleone for theorem and proposition (and the other statements), thmstyletwo for example and remark,
% thmstylethree for definition. These style names exist only in sn-jnl (amsthm falls back to plain with a warning
% under another class): a port to another class sets them back to definition, plain and remark.

% Numbering within sections
\numberwithin{equation}{section}

% One shared counter within a section: Definition 2.1, Theorem 2.2, Lemma 2.3, ...
\declaretheorem[name=Definition,  style=thmstylethree, numberwithin=section]{definition}
\declaretheorem[name=Theorem,     style=thmstyleone,   sibling=definition]{theorem}
\declaretheorem[name=Lemma,       style=thmstyleone,   sibling=definition]{lemma}
\declaretheorem[name=Claim,       style=thmstyleone,   sibling=definition]{claim}
\declaretheorem[name=Corollary,   style=thmstyleone,   sibling=definition]{corollary}
\declaretheorem[name=Proposition, style=thmstyleone,   sibling=definition]{proposition}
\declaretheorem[name=Remark,      style=thmstyletwo,   sibling=definition]{remark}
\declaretheorem[name=Example,     style=thmstyletwo,   sibling=definition]{example}

% "problem" environment: framed statement of a computational problem
\newtcolorbox{problem}[1][]{%
  enhanced,
  colback=white,      % background color
  colframe=black,     % border color
  boxrule=0.8pt,      % border thickness
  arc=2pt,            % corner roundness
  left=6pt, right=6pt, top=6pt, bottom=6pt,
  title=\textbf{Problem: #1},
}

% Algorithm headers: an algorithm states Input and Output.
% algpseudocode (loaded in main.tex, before this file) would print "Require:" and "Ensure:".
\algrenewcommand\algorithmicrequire{\textbf{Input:}}
\algrenewcommand\algorithmicensure{\textbf{Output:}}
